\section*{अध्याय \ref{ch:g11:deriv} --- अवकलन: पहला पाठ्यक्रम}

\begin{solution}{exo:g11:deriv:1}
$a = 2$ पर $f(x) = x^2$ के लिए:
\[
\frac{(2+h)^2 - 4}{h} = \frac{4h + h^2}{h} = 4 + h
\xrightarrow[h\to0]{} 4,
\]
इसलिए $f'(2) = 4$। $a = 1$ पर $g(x) = x^2 + 3x$ के लिए: $g(1) = 4$ और
\[
\frac{(1+h)^2 + 3(1+h) - 4}{h} = \frac{5h + h^2}{h} = 5 + h
\xrightarrow[h\to0]{} 5,
\]
इसलिए $g'(1) = 5$।
\end{solution}

\begin{solution}{exo:g11:deriv:2}
$f'(x) = 12x^2 - 4x + 1$।

$g'(x) = \frac{5}{2\sqrt x} - \frac{3}{x^2}$ ($\intoo{0}{+\infty}$ पर)।

$h'(x) = 2(x^2 - 3) + (2x+1)(2x) = 2x^2 - 6 + 4x^2 + 2x
= 6x^2 + 2x - 6$ (गुणनफल का नियम)।
\end{solution}

\begin{solution}{exo:g11:deriv:3}
$f(x) = x^2 - 3x$, $a = 2$: $f(2) = -2$, $f'(x) = 2x - 3$, इसलिए $f'(2) = 1$; \omterm{def:g11:deriv:tangent}{स्पर्श रेखा} $y = -2 + (x - 2) = x - 4$।

$f(x) = \frac1x$, $a = 1$: $f(1) = 1$, $f'(1) = -1$; \omterm{def:g11:deriv:tangent}{स्पर्श रेखा} $y = 1 - (x - 1) = 2 - x$।

$f(x) = \sqrt x$, $a = 4$: $f(4) = 2$, $f'(4) = \frac{1}{2\sqrt4} =
\frac14$; \omterm{def:g11:deriv:tangent}{स्पर्श रेखा} $y = 2 + \frac14(x - 4) = \frac{x}{4} + 1$।
\end{solution}

\begin{solution}{exo:g11:deriv:4}
$f'(x) = \frac{(x-1) - (x+2)}{(x-1)^2} = \frac{-3}{(x-1)^2}$।

$g'(x) = \frac{2x(x^2+1) - x^2 \times 2x}{(x^2+1)^2}
= \frac{2x}{(x^2+1)^2}$।

$h'(x) = -\frac{2x+1}{(x^2+x+1)^2}$ ($u = 1$ के साथ भागफल का नियम)।
\end{solution}

\begin{solution}{exo:g11:deriv:5}
$f'(x) = 3x^2 - 12x + 9 = 3(x^2 - 4x + 3) = 3(x-1)(x-3)$: $\intcc{1}{3}$ के बाहर धनात्मक, भीतर ऋणात्मक। इसलिए $f$ $\intoc{-\infty}{1}$ पर \omterm{def:g11:func:monotone}{वर्धमान}, $\intcc{1}{3}$
पर \omterm{def:g10:functions:variations}{ह्रासमान} और $\intco{3}{+\infty}$ पर \omterm{def:g11:func:monotone}{वर्धमान} है, जिसमें स्थानीय \omterm{def:g10:functions:extrema}{अधिकतम} $f(1) = 1 - 6 + 9 - 2 = 2$ और स्थानीय
\omterm{def:g10:functions:extrema}{न्यूनतम} $f(3) = 27 - 54 + 27 - 2 = -2$ है।
\end{solution}

\begin{solution}{exo:g11:deriv:6}
$\intoo{-1}{+\infty}$ पर:
\[
f'(x) = \frac{2x(x+1) - (x^2+3)}{(x+1)^2}
= \frac{x^2 + 2x - 3}{(x+1)^2}
= \frac{(x+3)(x-1)}{(x+1)^2}.
\]
$x > -1$ के लिए $x + 3$ और $(x+1)^2$ दोनों धनात्मक हैं, इसलिए $f'(x)$ का चिह्न $x - 1$ जैसा
है: $f$ $\intoc{-1}{1}$ पर \omterm{def:g10:functions:variations}{ह्रासमान} और $\intco{1}{+\infty}$ पर \omterm{def:g11:func:monotone}{वर्धमान} है, और उसका \omterm{def:g10:functions:extrema}{न्यूनतम} $f(1) = \frac{4}{2} = 2$ है।
\end{solution}

\begin{solution}{exo:g11:deriv:7}
$f'(x) = 2x - 4$।

\emph{1.} क्षैतिज \omterm{def:g11:deriv:tangent}{स्पर्श रेखा}: $x = 2$ पर $f'(x) = 0$; बिंदु $(2, f(2)) = (2, 1)$ है (\omterm{def:g11:quad:parabola}{परवलय} का शीर्ष)।

\emph{2.} $y = 2x + 1$ के समांतर होने का अर्थ है \omterm{def:g10:reffunc:affine}{प्रवणता} $2$: $2x - 4 = 2$, इसलिए $x = 3$,
जिससे बिंदु $(3, f(3)) = (3, 2)$ मिलता है।
\end{solution}

\begin{solution}{exo:g11:deriv:8}
चौड़ाई $x$ (बाड़ वाली दो भुजाएँ) और लंबाई $60 - 2x$ के साथ: $\intoo{0}{30}$ पर $A(x) = x(60 - 2x) = 60x - 2x^2$। तब
$A'(x) = 60 - 4x$, जो $x = 15$ से पहले धनात्मक और बाद में ऋणात्मक है: क्षेत्रफल $x = 15$ पर
\omterm{def:g10:functions:extrema}{अधिकतम} होता है, और $15 \times 30$ m का बाड़ा $450$ m$^2$ क्षेत्रफल देता है। \cref{ch:g11:quad} का
शीर्ष-सूत्र $\alpha = -\frac{60}{2 \times (-2)} = 15$ वही उत्तर कलन के बिना दे देता है।
\end{solution}

\begin{solution}{exo:g11:deriv:9}
मान लीजिए $\intoo{0}{+\infty}$ पर $f(x) = \frac{x+1}{2} - \sqrt x$ है। तब
\[
f'(x) = \frac12 - \frac{1}{2\sqrt x} = \frac{\sqrt x - 1}{2\sqrt x},
\]
जिसका चिह्न $\sqrt x - 1$ जैसा है: $\intoo{0}{1}$ पर ऋणात्मक, $\intoo{1}{+\infty}$ पर धनात्मक। इसलिए $f$ पहले
घटता है फिर बढ़ता है, और उसका \omterm{def:g10:functions:extrema}{न्यूनतम} $f(1) = 1 - 1 = 0$ है। अतः सभी $x > 0$ के लिए $f(x) \geq 0$,
अर्थात् $\sqrt x \leq \frac{x+1}{2}$, और समता केवल $x = 1$ पर।
\end{solution}

\begin{solution}{exo:g11:deriv:10}
$\pi r^2 h = 500$ से $h = \frac{500}{\pi r^2}$, इसलिए
\[
S(r) = 2\pi r^2 + 2\pi r \cdot \frac{500}{\pi r^2}
= 2\pi r^2 + \frac{1000}{r},
\qquad
S'(r) = 4\pi r - \frac{1000}{r^2}.
\]
$4\pi r^3 = 1000$ होने पर $S'(r) = 0$, अर्थात् $r^3 = \frac{250}{\pi}$, इसलिए $r = \sqrt[3]{250/\pi}$ ($\approx 4.3$ cm)। चूँकि $S'(r) =
\frac{4\pi r^3 - 1000}{r^2}$ इस
मान से पहले ऋणात्मक और बाद में धनात्मक है, यह \omterm{def:g10:functions:extrema}{न्यूनतम} है।
\end{solution}

\begin{solution}{exo:g11:deriv:11}
\omterm{def:g10:coordgeom:system}{भुज} $a$ वाले बिंदु पर $y = x^2$ की \omterm{def:g11:deriv:tangent}{स्पर्श रेखा} $y = a^2 + 2a(x - a) = 2ax - a^2$ है। वह $P(1, -3)$ से तब होकर
जाती है जब $-3 = 2a - a^2$ हो, अर्थात् $a^2 - 2a - 3 = 0$, जिसके \omterm{def:g11:quad:discriminant}{मूल} $a = 3$ और $a = -1$ हैं ($\Delta = 16$)।
इसलिए $P$ से होकर \emph{दो} \omterm{def:g11:deriv:tangent}{स्पर्श रेखाएँ} जाती हैं:
\[
y = 6x - 9 \quad (a = 3),
\qquad
y = -2x - 1 \quad (a = -1).
\]
\end{solution}

\begin{solution}{exo:g11:deriv:12}
\cref{ex:g11:deriv:study} से $f$ स्थानीय \omterm{def:g10:functions:extrema}{अधिकतम} $f(-1) = 3$ तक बढ़ता है, फिर स्थानीय \omterm{def:g10:functions:extrema}{न्यूनतम} $f(1) = -1$
तक घटता है, और फिर बढ़ता है; और बहुत बड़े धनात्मक (क्रमशः ऋणात्मक) $x$ के
लिए $x^3$ हावी हो जाता है और $f$ मनमाने बड़े (क्रमशः छोटे) मान लेता है।
\omterm{def:g10:functions:table}{परिवर्तन-सारणी} पर क्षैतिज रेखाएँ $y = k$ आरपार पढ़ने पर:
\begin{itemize}
  \item $k > 3$ या $k < -1$: रेखा वक्र से एक बार मिलती है --- $1$ हल;
  \item $k = 3$ या $k = -1$: रेखा किसी मोड़-बिंदु को छूती है और एक दूसरी शाखा
        को पार करती है --- $2$ हल;
  \item $-1 < k < 3$: रेखा तीनों शाखाओं को पार करती है --- $3$ हल।
\end{itemize}
\end{solution}

\begin{solution}{pb:g11:deriv:1}
\textbf{1.} $f'(x) = 6x - 5$; $g'(x) = 3x^2 - 12$। परिभाषा से: $\frac{(1 + h)^2 - 1}{h} = 2 + h \to 2$: $1$ पर $x^2$ का \omterm{def:g11:deriv:derivative}{अवकलज} $2$
है।

\textbf{2.} $g(2) = 8 - 24 = -16$ और $g'(2) = 12 - 12 = 0$: \omterm{def:g11:deriv:tangent}{स्पर्श रेखा} क्षैतिज रेखा $y = -16$ है। क्षैतिज
\omterm{def:g11:deriv:tangent}{स्पर्श रेखा} किसी क्रांतिक बिंदु का संकेत है --- यहाँ एक स्थानीय \omterm{def:g10:functions:extrema}{न्यूनतम},
जैसा प्रश्न 5 की \omterm{def:g10:functions:table}{परिवर्तन-सारणी} पुष्ट करती है।

\textbf{3.} $(a, a^2)$ पर \omterm{def:g10:reffunc:affine}{प्रवणता} $2a$: $y = a^2 + 2a(x - a) = 2ax - a^2$।

\textbf{4.} $x = 0$ पर: $y = -a^2$: \omterm{def:g11:deriv:tangent}{स्पर्श रेखा} अक्ष को मूल बिंदु से $a^2$ गहराई पर
काटती है --- ठीक उतना ही \emph{नीचे} जितना $P$ ऊपर बैठा है। नुस्ख़ा: ऊँचाई
$h$ वाले बिंदु पर \omterm{def:g11:deriv:tangent}{स्पर्श रेखा} खींचने के लिए उसे मूल बिंदु से $h$ गहराई
वाले अक्ष-बिंदु से जोड़ दीजिए। (रेखाचित्र-पटल पर किसी \omterm{def:g11:deriv:derivative}{अवकलज} की आवश्यकता
नहीं।)

\textbf{5.} $g'(x) = 3(x^2 - 4)$: $\intoo{-\infty}{-2}$ पर धनात्मक, $\intoo{-2}{2}$ पर ऋणात्मक, और उसके आगे धनात्मक।
$g$ बढ़कर स्थानीय \omterm{def:g10:functions:extrema}{अधिकतम} $g(-2) = 16$ तक जाता है, गिरकर स्थानीय \omterm{def:g10:functions:extrema}{न्यूनतम} $g(2) = -16$ तक,
फिर बढ़ता है।

\textbf{6.} $F\left(0, \frac14\right)$ और $T(0, -a^2)$: $FT = \frac14 + a^2$।

\textbf{7.} $FP = a^2 + \frac14 = FT$: $F$ पर समद्विबाहु, इसलिए आधार के दोनों कोण बराबर हैं:
$\widehat{FTP} = \widehat{FPT}$।

\textbf{8.} $P$ पर आती हुई किरण ऊर्ध्वाधर है, अतः रेखा $(TF)$ ($y$-अक्ष)
के समांतर। \omterm{def:g11:deriv:tangent}{स्पर्श रेखा} $(TP)$ दोनों समांतर रेखाओं को काटती है, इसलिए किरण और
\omterm{def:g11:deriv:tangent}{स्पर्श रेखा} के बीच का कोण $\widehat{FTP}$ के बराबर है (एकांतर कोण) --- और वह $\widehat{FPT}$ के
बराबर है, जो $[PF]$ और \omterm{def:g11:deriv:tangent}{स्पर्श रेखा} के बीच का कोण है। $P$ पर \omterm{def:g11:deriv:tangent}{स्पर्श रेखा} के
दोनों ओर बराबर कोण।

\textbf{9.} परावर्तन के नियम से $P$ पर ऊर्ध्वाधर आती किरण दूसरी ओर उसी
बराबर कोण वाली रेखा के अनुदिश लौटती है --- और प्रश्न 8 से वह $[PF]$ के अनुदिश
है: \omterm{pb:g11:quad:1}{नाभि} से होकर, चाहे $a$ कुछ भी हो। उलटी दिशा में, $F$ पर जन्मा प्रकाश
दर्पण के हर बिंदु से अक्ष के समांतर निकलता है। तश्तरियाँ बटोरती हैं,
हेडलाइट पुंज फेंकती हैं: सिद्ध।

\textbf{10.} $a = 1$: \omterm{def:g11:deriv:tangent}{स्पर्श रेखा} $y = 2x - 1$, $T(0, -1)$; $FP = \sqrt{1 + \left(1 - \frac14\right)^2} =
\sqrt{\frac{25}{16}} = \frac54$; $FT = \frac14 + 1 = \frac54$: समद्विबाहु,
जैसा वादा था।

\textbf{11.} $x_n$ पर \omterm{def:g11:deriv:tangent}{स्पर्श रेखा} $y = f(x_n) + f'(x_n)(x - x_n)$ है; वह $y = 0$ को वहाँ काटती है जहाँ
$x = x_n - \frac{f(x_n)}{f'(x_n)}$।

\textbf{12.} $x_1 = 1 - \frac{-1}{2} = \frac32$; $x_2 = \frac32 - \frac{1/4}{3} = \frac{17}{12}$; $x_3 = \frac{577}{408}$: $\sqrt2$ के हीरोन-सन्निकटन, अर्थात् \cref{pb:g10:functions:1} की
स्थिर-बिंदु वाली मशीन। सर्वसमिका: $x - \frac{x^2 - 2}{2x} = \frac{2x^2 - x^2 + 2}{2x}
= \frac{x^2 + 2}{2x} = \frac12\left(x + \frac2x\right)$ --- हीरोन का नुस्ख़ा वस्तुतः $x^2 - 2$
पर लगाई गई न्यूटन की विधि \emph{ही} है, दो सहस्राब्दी पहले।

\textbf{13.} $x_1 = 5 - \frac{125 - 100}{75} = \frac{14}{3}
\approx 4.6667$; $x_2 \approx 4.6667 - \frac{1.63}{65.33}
\approx 4.6417$। कैलकुलेटर: $\sqrt[3]{100} = 4.6416\ldots$ --- दो क़दमों में चार सही दशमलव
स्थान।

\textbf{14.} $f'(2) = 0$: आरंभिक बिंदु पर \omterm{def:g11:deriv:tangent}{स्पर्श रेखा} क्षैतिज है (प्रश्न 2) और वह
अक्ष को कभी नहीं काटती --- सूत्र शून्य से भाग दे देता है। चेतावनी: न्यूटन को
क्रांतिक बिंदुओं से दूर (और, अधिक सामान्य रूप से, अभीष्ट \omterm{def:g11:quad:discriminant}{मूल} के पर्याप्त
पास) से शुरू कीजिए।

\textbf{15.} द्विभाजन हर क़दम पर एक द्विआधारी अंक जोड़ता है; न्यूटन हर क़दम
पर सही अंकों की संख्या मोटे तौर पर \emph{दुगुनी} कर देता है (यहाँ $1 \to 3 \to 6$)।
जब हर क़दम महँगा हो --- लाखों प्राचल, वास्तविक-समय नौवहन --- तब दुगुना करना
घिसटने से बेहतर है, और इसीलिए इतने सारे सिलिकॉन के भीतर \omterm{def:g11:deriv:tangent}{स्पर्श रेखा} दौड़
रही है।

\textbf{16.} $d^2 = 900 - w^2$, इसलिए $0 < w < 30$ के लिए $S(w) = w(900 - w^2) = 900w - w^3$।

\textbf{17.} $w = \sqrt{300} = 10\sqrt3 \approx 17.3$ cm पर $S'(w) = 900 - 3w^2 = 0$ ($S' > 0$ पहले, $< 0$ बाद में: एक \omterm{def:g10:functions:extrema}{अधिकतम})। फिर
$d = \sqrt{600} = 10\sqrt6 \approx 24.5$ cm।

\textbf{18.} $\frac{d^2}{w^2} = \frac{600}{300} = 2$, इसलिए $d = w\sqrt2$: गहराई चौड़ाई का $\sqrt2$ गुना है --- वही
बढ़इयों का अनुपात (और व्यास को तीन भागों में बाँटने वाली रचना उसे परकार से
ही दे देती है, मौक़े पर कोई बीजगणित नहीं)।

\textbf{19.} वर्गाकार शहतीर: $w = d = \sqrt{450} \approx 21.2$ cm, दृढ़ता $w^3 = 450^{3/2} \approx 9\,546$। इष्टतम शहतीर: $S = 10\sqrt3 \times 600 = 6\,000\sqrt3 \approx 10\,392$।
लाभ: उसी लट्ठे से लगभग $8.9\,\%$ अधिक दृढ़ता --- \omterm{def:g11:deriv:derivative}{अवकलज} आरा-मिल का ख़र्च निकाल
देता है।

\textbf{20.} ज्यामिति: \omterm{def:g11:deriv:tangent}{स्पर्श रेखा} के बराबर कोणों ने एक परवलयज को
प्रकाश-संग्राहक बना दिया। संख्यात्मक गणना: स्पर्श रेखाओं पर फिसलते जाना
\omterm{def:g10:algebra:equation}{समीकरणों} को अंक-दुगुने करने की गति से हल कर देता है। इष्टतमीकरण: लुप्त होते
\omterm{def:g11:deriv:derivative}{अवकलज} ने लट्ठे का सबसे मज़बूत आयत ढूँढ़ निकाला। एक ही वस्तु, तीन पेशे ---
और कक्षा 12 की उत्तलता इसके ऊपर यह भी बता देगी कि हर \omterm{def:g11:deriv:tangent}{स्पर्श रेखा} के किस ओर
वक्र पड़ा है।
\end{solution}
